\documentclass[11pt,a4paper]{article}
\newcommand{\intestazione}{Trimero ad anello --- Schema dei quattro passi}
% =====================================================================
%  Preambolo comune ai documenti sul dimero (serie "che cosa abbiamo fatto")
%  Incluso con % =====================================================================
%  Preambolo comune ai documenti sul dimero (serie "che cosa abbiamo fatto")
%  Incluso con % =====================================================================
%  Preambolo comune ai documenti sul dimero (serie "che cosa abbiamo fatto")
%  Incluso con \input{preambolo_dimero} subito dopo \documentclass.
% =====================================================================
\usepackage[T1]{fontenc}
\usepackage[utf8]{inputenc}
\usepackage[main=italian,provide=*]{babel}
\usepackage{amsmath,amssymb,mathtools}
\usepackage{graphicx}
\usepackage{booktabs}
\usepackage{colortbl}
\usepackage{xcolor}
\usepackage{geometry}
\usepackage{placeins}
\usepackage{caption}
\usepackage{enumitem}
\usepackage{fancyhdr}
\usepackage{needspace}
\usepackage{tikz}
\usetikzlibrary{arrows.meta,positioning,shapes.geometric,calc,fit,backgrounds}
\usepackage{hyperref}
\hypersetup{colorlinks=true,linkcolor=blu,citecolor=blu,urlcolor=blu}

\geometry{margin=2.5cm}

% --- colori coerenti con quelli delle figure (palette Okabe-Ito)
\definecolor{blu}{HTML}{0072B2}
\definecolor{arancio}{HTML}{E69F00}
\definecolor{verdes}{HTML}{009E73}
\definecolor{rossos}{HTML}{D55E00}
\definecolor{violas}{HTML}{CC79A7}
\definecolor{grigetto}{HTML}{E8E8E8}

\captionsetup{font=small,labelfont=bf,width=0.92\textwidth}

\pagestyle{fancy}
\fancyhf{}
\fancyhead[L]{\small\itshape\intestazione}
\fancyhead[R]{\small\thepage}
\renewcommand{\headrulewidth}{0.3pt}

% --- notazione
\newcommand{\ket}[1]{\lvert #1\rangle}
\newcommand{\bra}[1]{\langle #1 \rvert}
\newcommand{\media}[1]{\langle #1 \rangle}
\newcommand{\vs}[1]{\mathbf{s}_{#1}}

% --- riquadro "in una riga": il messaggio da portare a casa
\newcommand{\messaggio}[1]{%
  \begin{center}
  \begin{minipage}{0.93\textwidth}
  \setlength{\fboxsep}{9pt}%
  \colorbox{grigetto}{\begin{minipage}{\dimexpr\textwidth-2\fboxsep}
  \small\textbf{In una riga.}\enspace #1
  \end{minipage}}
  \end{minipage}
  \end{center}
}

% --- riquadro "come lo sappiamo": la verifica numerica dietro un'affermazione
\newcommand{\verifica}[1]{%
  \begin{center}
  \begin{minipage}{0.93\textwidth}
  \small\textcolor{verdes}{\rule{2.2em}{1.1pt}}\enspace
  \textcolor{verdes}{\textbf{Come lo sappiamo.}}\enspace #1
  \end{minipage}
  \end{center}
  \medskip
}
 subito dopo \documentclass.
% =====================================================================
\usepackage[T1]{fontenc}
\usepackage[utf8]{inputenc}
\usepackage[main=italian,provide=*]{babel}
\usepackage{amsmath,amssymb,mathtools}
\usepackage{graphicx}
\usepackage{booktabs}
\usepackage{colortbl}
\usepackage{xcolor}
\usepackage{geometry}
\usepackage{placeins}
\usepackage{caption}
\usepackage{enumitem}
\usepackage{fancyhdr}
\usepackage{needspace}
\usepackage{tikz}
\usetikzlibrary{arrows.meta,positioning,shapes.geometric,calc,fit,backgrounds}
\usepackage{hyperref}
\hypersetup{colorlinks=true,linkcolor=blu,citecolor=blu,urlcolor=blu}

\geometry{margin=2.5cm}

% --- colori coerenti con quelli delle figure (palette Okabe-Ito)
\definecolor{blu}{HTML}{0072B2}
\definecolor{arancio}{HTML}{E69F00}
\definecolor{verdes}{HTML}{009E73}
\definecolor{rossos}{HTML}{D55E00}
\definecolor{violas}{HTML}{CC79A7}
\definecolor{grigetto}{HTML}{E8E8E8}

\captionsetup{font=small,labelfont=bf,width=0.92\textwidth}

\pagestyle{fancy}
\fancyhf{}
\fancyhead[L]{\small\itshape\intestazione}
\fancyhead[R]{\small\thepage}
\renewcommand{\headrulewidth}{0.3pt}

% --- notazione
\newcommand{\ket}[1]{\lvert #1\rangle}
\newcommand{\bra}[1]{\langle #1 \rvert}
\newcommand{\media}[1]{\langle #1 \rangle}
\newcommand{\vs}[1]{\mathbf{s}_{#1}}

% --- riquadro "in una riga": il messaggio da portare a casa
\newcommand{\messaggio}[1]{%
  \begin{center}
  \begin{minipage}{0.93\textwidth}
  \setlength{\fboxsep}{9pt}%
  \colorbox{grigetto}{\begin{minipage}{\dimexpr\textwidth-2\fboxsep}
  \small\textbf{In una riga.}\enspace #1
  \end{minipage}}
  \end{minipage}
  \end{center}
}

% --- riquadro "come lo sappiamo": la verifica numerica dietro un'affermazione
\newcommand{\verifica}[1]{%
  \begin{center}
  \begin{minipage}{0.93\textwidth}
  \small\textcolor{verdes}{\rule{2.2em}{1.1pt}}\enspace
  \textcolor{verdes}{\textbf{Come lo sappiamo.}}\enspace #1
  \end{minipage}
  \end{center}
  \medskip
}
 subito dopo \documentclass.
% =====================================================================
\usepackage[T1]{fontenc}
\usepackage[utf8]{inputenc}
\usepackage[main=italian,provide=*]{babel}
\usepackage{amsmath,amssymb,mathtools}
\usepackage{graphicx}
\usepackage{booktabs}
\usepackage{colortbl}
\usepackage{xcolor}
\usepackage{geometry}
\usepackage{placeins}
\usepackage{caption}
\usepackage{enumitem}
\usepackage{fancyhdr}
\usepackage{needspace}
\usepackage{tikz}
\usetikzlibrary{arrows.meta,positioning,shapes.geometric,calc,fit,backgrounds}
\usepackage{hyperref}
\hypersetup{colorlinks=true,linkcolor=blu,citecolor=blu,urlcolor=blu}

\geometry{margin=2.5cm}

% --- colori coerenti con quelli delle figure (palette Okabe-Ito)
\definecolor{blu}{HTML}{0072B2}
\definecolor{arancio}{HTML}{E69F00}
\definecolor{verdes}{HTML}{009E73}
\definecolor{rossos}{HTML}{D55E00}
\definecolor{violas}{HTML}{CC79A7}
\definecolor{grigetto}{HTML}{E8E8E8}

\captionsetup{font=small,labelfont=bf,width=0.92\textwidth}

\pagestyle{fancy}
\fancyhf{}
\fancyhead[L]{\small\itshape\intestazione}
\fancyhead[R]{\small\thepage}
\renewcommand{\headrulewidth}{0.3pt}

% --- notazione
\newcommand{\ket}[1]{\lvert #1\rangle}
\newcommand{\bra}[1]{\langle #1 \rvert}
\newcommand{\media}[1]{\langle #1 \rangle}
\newcommand{\vs}[1]{\mathbf{s}_{#1}}

% --- riquadro "in una riga": il messaggio da portare a casa
\newcommand{\messaggio}[1]{%
  \begin{center}
  \begin{minipage}{0.93\textwidth}
  \setlength{\fboxsep}{9pt}%
  \colorbox{grigetto}{\begin{minipage}{\dimexpr\textwidth-2\fboxsep}
  \small\textbf{In una riga.}\enspace #1
  \end{minipage}}
  \end{minipage}
  \end{center}
}

% --- riquadro "come lo sappiamo": la verifica numerica dietro un'affermazione
\newcommand{\verifica}[1]{%
  \begin{center}
  \begin{minipage}{0.93\textwidth}
  \small\textcolor{verdes}{\rule{2.2em}{1.1pt}}\enspace
  \textcolor{verdes}{\textbf{Come lo sappiamo.}}\enspace #1
  \end{minipage}
  \end{center}
  \medskip
}

\usepackage{pdflscape}
\usepackage{array}

\title{\bfseries Schema riassuntivo\\[2pt]
\large Punti di lavoro e collegamenti nei quattro passi, trimero ad anello}
\author{Samuele Schiavi \\ \small Tesi triennale in Fisica, Università di Parma
--- Relatore: Prof.\ Alessandro Chiesa}
\date{}

\begin{document}
\maketitle
\thispagestyle{fancy}

Questo documento non introduce fisica nuova: raccoglie in un unico posto,
per ciascuno dei quattro passi, il punto di lavoro, lo stato iniziale, il
risultato chiave e il collegamento con quanto viene dopo --- pensato per
essere consultato, non letto in ordine. Mirror diretto dello schema
analogo della serie sul dimero.

\section{Tabella d'insieme}

\renewcommand{\arraystretch}{1.35}
\begin{landscape}
\thispagestyle{empty}
\begin{center}
\small
\setlength{\tabcolsep}{4pt}
\begin{tabular}{@{}>{\bfseries\raggedright}p{2.7cm} p{5.15cm} p{5.15cm} p{5.15cm} p{5.15cm}@{}}
\toprule
\rowcolor{grigetto}
 & \textbf{Il sistema} & \textbf{Il VQE} & \textbf{La dinamica (Trotter)} & \textbf{I correlatori} \\
\midrule
Obiettivo &
Spettro esatto (Kambe); frustrazione; due opzioni DM, effetti opposti &
Trovare il fondamentale; confronto RBS contro $W$ sotto DM &
$e^{-iHt}$ con due livelli di Trotter interno (scambio e DM annidati sui
tre legami) &
Correlatori $C_{ij}^{\alpha\beta}(t)$, circuito a 4 qubit (3+ancilla) \\
\addlinespace
Punto/i di lavoro &
sweep in $b/J$; $D=0$ vs opzioni DM A/B\textsuperscript{*} ($D/J=0.15$) &
Scenario A\textsuperscript{**}: $b/J=b_c=2.4$, $D/J=0.15$ (opzione DM B)
--- risultati principali qui, ma verificato anche su sweep (13 punti,
$\mathcal F=1$ ovunque) &
$R_0$: $b/J=0.05$, $D/J=1.93$ --- \emph{dichiarato provvisorio} &
Scenario A\textsuperscript{**} --- stesso punto singolo del VQE \\
\addlinespace
Stato iniziale &
nessuno (solo diagonalizzazione) &
$X\ket{000}$ per costruzione (registro), poi porte parametrizzate &
registro $\ket{000}$, nessuna preparazione &
fondamentale esatto, fase fissata reale --- più un confronto dedicato con
la preparazione reale via $W$-2q.6 \\
\addlinespace
Risultato chiave &
incrocio a $b_c/J=2.4$ ($D{=}0$); a $D/J=0.15$, opzione DM A non lo apre
(gap $2.3\times10^{-8}$ a $b/J{=}2.372$), opzione DM B sì (gap vero
$0.254$ a $b/J{=}2.407$) --- entrambi i minimi spostati da $b_c$, non a
$b_c$ fisso &
RBS-2q ha un tetto vero ($\mathcal F=0.99991$, non si ripara con più giri);
$W$-2q.6 raggiunge $\mathcal F=1$ esatto, con meno gate &
errore $\sim1/N$ (operatori), $\sim1/N^2$ (fidelity), ma convergenza molto
più lenta del dimero: servono centinaia-migliaia di passi, non decine &
4 zeri strutturali su 81 (contro 32 nulli solo a $t{=}0$); ricchezza
spettrale indipendente dall'ampiezza, come sul dimero; preparazione VQE
reale contro esatta: scarto medio $2.9\times10^{-7}$, nessuna correlazione
con l'ampiezza (a differenza del dimero, l'ansatz qui è quasi esatto) \\
\addlinespace
Collegamento &
motiva la scelta dell'opzione DM B per tutto il resto della serie &
fornisce l'ansatz canonico ($W$-2q.6) riusato nei due passi seguenti, ad
ogni punto proprio &
introduce il costo in gate (15 CNOT/passo) che la pipeline di rumore
riprende direttamente &
chiude la serie; circuito più lungo del pacchetto, primo banco di prova
per il sistema aperto \\
\bottomrule
\end{tabular}

\bigskip
\noindent{\footnotesize
\textsuperscript{*}\textbf{Opzioni DM A/B} (Documento 1): due forme
candidate del termine di Dzyaloshinskii--Moriya, non due punti di lavoro
--- l'Opzione A non apre mai l'incrocio (commuta con $S_{12}^2$), l'Opzione
B lo apre (non commuta). Per tutta la serie a valle del Documento 1 si usa
solo l'Opzione B.\quad
\textsuperscript{**}\textbf{Scenario A} (Documento 2 in poi): un punto di
lavoro del VQE/correlatori, $b/J=b_c=2.4$, $D/J=0.15$ --- \emph{non} legato
all'Opzione A del DM (che qui non è usata: Scenario A lavora sempre con
l'Opzione B). La stessa lettera ``A'' indica due cose diverse in due
contesti diversi; tenerle distinte.}
\end{center}
\end{landscape}

\clearpage

\section{Il sistema}

Diagonalizzazione esatta della matrice $8\times8$, decomposizione di Kambe
in tre multipletti (A, B, C). Il risultato che conta per il resto del
lavoro: la frustrazione geometrica (ciclo dispari) è strutturale, mentre
l'apertura dell'incrocio dipende da \emph{quale} forma di DM si accende.
Non è una scelta libera fra due candidati equivalenti: verificato
sistematicamente (tutte le combinazioni di segno sui tre legami, contro la
simmetria di scambio $1\leftrightarrow2$) che l'Opzione A è l'\textbf{unica}
forma che rispetta quella simmetria; l'Opzione B è la proposta del
relatore, più semplice ma non derivata dalla simmetria --- e l'unica delle
due che apre davvero l'incrocio.

\messaggio{Con un solo legame il dimero non poteva porre la domanda ``quale
DM?''. Con tre legami la risposta cambia la fisica: solo l'Opzione B apre
l'incrocio, ed è per questo l'unica scelta usata nel resto della serie.}

\section{Il VQE}

Due famiglie di ansatz a confronto, stessa struttura esterna (sei
parametri: tre sui legami, tre rotazioni $R_y$ finali), diverso blocco a
due qubit (RBS o $W$). Senza DM, equivalenti ($\mathcal F=1$ per entrambe).
Con DM acceso, a Scenario A: RBS ha un tetto reale ($\mathcal
F=0.9999110739$), verificato non rimuovibile aggiungendo altri giri dello
stesso blocco; $W$ raggiunge l'esatto ($\mathcal F=1$ a sedici cifre), con
meno gate a un qubit. Ansatz canonico per il seguito: \textbf{$W$-2q.6}.

Verificato anche oltre il punto singolo: multistart su 13 valori di $b/J$
fra $0.3$ e $4.5$ (DM Opzione B acceso), incluso il punto dell'incrocio
stesso --- $\mathcal F=1.00000000$ ovunque, nessun ramo di preparazione
necessario.

\messaggio{Sul dimero, RBS batte $W$ (stesso potere espressivo, meno gate).
Sul trimero, per questo particolare stato di partenza, $W$ batte RBS anche
in potere espressivo, non solo in costo --- un'inversione rispetto al
dimero che vale la pena tenere a mente quando si generalizza da un sistema
all'altro.}

\section{La dinamica}

Struttura a tre livelli (non due, come sul dimero): campo e scambio isotropo
fattorizzano esattamente fra loro, ma lo scambio stesso richiede un Trotter
interno sui tre legami condivisi --- stesso per il termine DM. Verificato:
circuito Qiskit e matrice compatta coincidono a precisione macchina
($4\times10^{-15}$) --- controllo di coerenza interna (la matrice compatta
usa l'operatore dello stesso circuito, non una derivazione indipendente),
a patto che rispetti anch'essa il Trotter interno (non basta
l'esponenziale diretto di $H_0$). Convergenza
molto più lenta del dimero: centinaia-migliaia di passi contro decine, per
via dei due livelli interni aggiuntivi. Costo: 15 CNOT per passo dopo
compilazione ottimizzata (dimezzato rispetto alla traduzione diretta, 30).

\messaggio{Il Trotter interno sui legami condivisi non è un dettaglio
implementativo: è la ragione fisica per cui il trimero converge molto più
lentamente del dimero a parità di soglia di accuratezza --- un costo che
il dimero, con un solo legame, non paga mai.}

\section{I correlatori}

Stesso schema Hadamard test del dimero, esteso a un registro di tre qubit
più ancilla (4 qubit totali; ancilla = qubit 3, non 0 come sul dimero).
Punto di lavoro Scenario A, lo stesso del VQE. Su 81 combinazioni: 32 nulle
solo all'istante iniziale (stesso argomento del dimero, stato reale), ma
solo \textbf{4 strutturalmente nulle per ogni $t$} --- tutte sullo stesso
sito (sito 3), non ancora spiegate in forma chiusa.

Chiuso anche il confronto con la preparazione reale (Documento 4, ultima
sezione): a differenza del dimero, dove l'ansatz base imperfetto
($\mathcal F=0.940$) produce uno scarto marcato e dipendente dal
correlatore, qui l'ansatz canonico $W$-2q.6 è quasi esatto
($1-\mathcal F\sim4\times10^{-13}$) --- lo scarto propagato resta
uniformemente minuscolo ($\sim3\times10^{-7}$), senza la stessa struttura
ricco/monocromatico osservata sul dimero: non c'è un ansatz imperfetto da
confrontare, quindi non c'è la stessa storia da raccontare.

\messaggio{Il numero di combinazioni cresce da 36 (dimero) a 81 (trimero),
ma la disciplina resta identica: distinguere zeri iniziali da zeri
strutturali, e non confondere ampiezza grande con ricchezza spettrale ---
due criteri indipendenti, verificato di nuovo su un sistema più grande.}

\end{document}
